\documentclass{ximera}
%% Macros shared by every activity. Each activity loads this file with
%% %% Macros shared by every activity. Each activity loads this file with
%% %% Macros shared by every activity. Each activity loads this file with
%% \input{../preamble.tex}, and the xourse file with \input{preamble.tex}.

%% MathML intent for screen readers: \xmunits, \xmabs, \xmpoint and the
%% other commands of TailorSwiftBot.sty. The file is latex/xmintent.sty.
\usepackage{xmintent}

%% The interactive pieces of Ximera: \answer, multipleChoice and \choice,
%% selectAll, hint and the rest. Keep this line: without it those commands
%% are undefined and come out as plain text. The file is
%% latex/ximera-components.sty.
\usepackage{ximera-components}

%% Tables with the tabularray package. The settings rowhead, xmcolhead,
%% caption and summary make their header cells known to screen readers, on
%% the web page and in the tagged PDF.
\usepackage{tabularray}

\newcommand{\RR}{\mathbb{R}}
\newcommand{\dfn}{\textit}
, and the xourse file with \input{preamble.tex}.

%% MathML intent for screen readers: \xmunits, \xmabs, \xmpoint and the
%% other commands of TailorSwiftBot.sty. The file is latex/xmintent.sty.
\usepackage{xmintent}

%% The interactive pieces of Ximera: \answer, multipleChoice and \choice,
%% selectAll, hint and the rest. Keep this line: without it those commands
%% are undefined and come out as plain text. The file is
%% latex/ximera-components.sty.
\usepackage{ximera-components}

%% Tables with the tabularray package. The settings rowhead, xmcolhead,
%% caption and summary make their header cells known to screen readers, on
%% the web page and in the tagged PDF.
\usepackage{tabularray}

\newcommand{\RR}{\mathbb{R}}
\newcommand{\dfn}{\textit}
, and the xourse file with \input{preamble.tex}.

%% MathML intent for screen readers: \xmunits, \xmabs, \xmpoint and the
%% other commands of TailorSwiftBot.sty. The file is latex/xmintent.sty.
\usepackage{xmintent}

%% The interactive pieces of Ximera: \answer, multipleChoice and \choice,
%% selectAll, hint and the rest. Keep this line: without it those commands
%% are undefined and come out as plain text. The file is
%% latex/ximera-components.sty.
\usepackage{ximera-components}

%% Tables with the tabularray package. The settings rowhead, xmcolhead,
%% caption and summary make their header cells known to screen readers, on
%% the web page and in the tagged PDF.
\usepackage{tabularray}

\newcommand{\RR}{\mathbb{R}}
\newcommand{\dfn}{\textit}


\title{Links Within a Page}

\begin{document}
\begin{abstract}
Cross-references and citations: what each link says and where it takes you.
\end{abstract}
\maketitle

This page exists to be tested. Each link in the last section should say a
number (never the name of a label such as ``eq:pythagoras''), and following
it should move you, and the keyboard focus, to what it names. With a screen
reader, press Enter on a link and listen: the reader should go on from the
equation, heading, theorem, problem or entry of the bibliography it names.

\section{Equations}\label{sec:equations}

A single numbered equation:
\begin{equation}\label{eq:pythagoras}
a^2+b^2=c^2
\end{equation}
This sentence follows the equation in the same paragraph, and should line
up with the rest of the text.

An \texttt{align} with a label on each row:
\begin{align}
x + y &= 3 \label{eq:sum}\\
x - y &= 1 \label{eq:difference}
\end{align}
and a display without a number:
\[ e^{i\pi} + 1 = 0 \]
followed by more of the same paragraph.

\subsection{Statements}\label{sec:statements}

\begin{theorem}\label{thm:parity}
The sum of two even numbers is even.
\end{theorem}

\begin{lemma}[Halving]\label{lem:halving}
Every even number is twice a whole number.
\end{lemma}

\begin{definition}\label{def:even}
A whole number is even if it is divisible by $2$.
\end{definition}

\begin{theorem}\label{thm:second}
The product of an even number and any whole number is even.
\end{theorem}

\section{Problems}\label{sec:problems}

\begin{problem}\label{prob:add}
What is $2+2$? $\answer{4}$
\end{problem}

\begin{problem}\label{prob:solve}
Using equations \eqref{eq:sum} and \eqref{eq:difference}, find $x$.
$x = \answer{2}$
\end{problem}

\section{Links}

\begin{itemize}
\item Equations: \eqref{eq:pythagoras}, \eqref{eq:sum} and
  \eqref{eq:difference}; without parentheses, \ref{eq:pythagoras}.
\item Sections: Section~\ref{sec:equations}, Subsection~\ref{sec:statements}
  and Section~\ref{sec:problems}.
\item Statements: Theorem~\ref{thm:parity}, Lemma~\ref{lem:halving},
  Definition~\ref{def:even} and Theorem~\ref{thm:second}.
\item Problems: Problem~\ref{prob:add} and Problem~\ref{prob:solve}.
\item Citations: \cite{Szego}; with a note, \cite[Theorem 8.1.1]{Szego};
  two at once, \cite{Szego,Knuth}.
\end{itemize}

\begin{thebibliography}{9}
\bibitem{Knuth} D. E. Knuth, \emph{The \TeX book}, Addison--Wesley, 1984.
\bibitem{Szego} G. Szeg\H{o}, \emph{Orthogonal Polynomials}, 4th edition,
  American Mathematical Society, 1975.
\end{thebibliography}
\end{document}
